\chapter*{Заключение}
\addcontentsline{toc}{chapter}{Заключение}
\label{ch:conclusion}

Все задачи второго этапа, предусмотренные п.~5.2.1 технического задания, решены полностью. Соответствие полученных результатов перечню результатов НИР (п.~3.1 ТЗ) и техническим требованиям (п.~4.1 ТЗ) приведено ниже; подробные сводки -- в таблицах~\ref{tab:tz-surf} и~\ref{tab:tz-light}.

Обзоры литературы (пп.~3.1.7 и 3.1.10 ТЗ):
\begin{itemize}
    \item проведён аналитический обзор научной литературы на тему эффективных методов дифференцируемого рендеринга для реконструкции или оптимизации конфигурации освещения (раздел~\ref{ch:light-review}): 21 ссылка, из них 15 -- на научные работы 2022--2026~гг.; в качестве аналога для сравнения выбрана система Mitsuba~3 с методом path replay backpropagation (пп.~4.1.15, 4.1.16 ТЗ);
    \item проведён аналитический обзор научной литературы на тему эффективных методов дифференцируемого рендеринга поверхностей (раздел~\ref{ch:chap1}): 23 ссылки, из них 15 -- на работы 2022--2026~гг.; в качестве аналога выбран метод релаксированной границы (Relaxed Boundary), дополнительно проведено сравнение с методом Vicini et al. (пп.~4.1.20, 4.1.21 ТЗ).
\end{itemize}

Научно-технические результаты.

Методы:
\begin{itemize}
    \item метод дифференцируемого рендеринга поверхностей, заданных разреженным множеством бриков функции расстояния со знаком (п.~3.1.11 ТЗ, раздел~\ref{ch:chap2}): реконструкция поверхности по изображениям с оценкой градиентов методом релаксированной границы, многоэпохальным повышением разрешения и разреживанием представления, при котором размер модели растёт приблизительно в 4 раза за эпоху вместо 8 раз у плотной сетки;
    \item метод дифференцируемого рендеринга для реконструкции конфигурации освещения (п.~3.1.8 ТЗ, раздел~\ref{ch:light-method}): реконструкция гониометрической диаграммы направленности и цвета источника, а также пространственно неоднородной текстуры излучения площадного источника по одному HDR-изображению сцены с учётом глобального освещения, основанная на однократном построении кэша световых путей и аналитическом вычислении градиента проходами по кэшу.
\end{itemize}

Модели:
\begin{itemize}
    \item модель излучения пространственно неоднородного площадного источника, задающая в каждой ячейке текстуры интенсивность, нормированный цвет и угловую ширину диаграммы направленности (раздел~\ref{ch:light-method});
    \item табличная модель излучения источника с гониометрической диаграммой направленности, совместимая с форматом IES LM-63 (раздел~\ref{ch:light-method}).
\end{itemize}

Алгоритмы:
\begin{itemize}
    \item алгоритм поиска пересечений луча и точек релаксированной границы внутри вокселя разреженного множества бриков методом Ньютона и аналитическим методом (раздел~\ref{ch:chap2});
    \item алгоритмы глобального и локального редистансинга разреженного множества бриков (раздел~\ref{ch:chap2});
    \item алгоритм вычисления изображения и градиента по кэшу световых путей на центральном и графическом процессорах без атомарных операций (раздел~\ref{ch:light-method});
    \item алгоритм решения задачи реконструкции текстуры излучения методом наименьших квадратов при фиксированной угловой модели (раздел~\ref{ch:light-method});
    \item алгоритм достраивания ненаблюдаемых направлений диаграммы направленности по симметриям, подтверждённым наблюдаемыми данными (раздел~\ref{ch:light-method}).
\end{itemize}

Программное обеспечение:
\begin{itemize}
    \item система LightCadRenderCore (п.~3.1.13 ТЗ, подраздел~\ref{sec:lightcad}) -- консольное приложение, выполняющее реконструкцию источника освещения по HDR-фотографии сцены, экспорт результата в файл IES, расчёт изображения сцены с восстановленным источником и сравнение с фотографией;
    \item программная реализация метода дифференцируемого рендеринга поверхностей на C++ с использованием Vulkan; вычислительные шейдеры на языке Slang порождаются транслятором kernel\_slicer (п.~4.1.1 ТЗ).
\end{itemize}

Иные научно-технические результаты: наборы сценариев испытаний, воспроизводящие сравнения с аналогами одним запуском (разделы~\ref{ch:chap4} и~\ref{chap:light-comparison}).

Результаты сравнения с аналогами и выполнение технических требований:
\begin{itemize}
    \item метод дифференцируемого рендеринга поверхностей (п.~3.1.12 ТЗ, раздел~\ref{ch:chap4}): PSNR реконструкции 12 моделей -- в среднем 38,3~дБ, не ниже 30,1~дБ на всех входных изображениях при требовании 25~дБ (п.~4.1.22 ТЗ); по сравнению с методом Relaxed Boundary на графическом процессоре RTX~2080 итерация реконструкции выполняется в 6,8--10,3 раза быстрее, все 7 эпох -- в среднем за 68~мин против 237~мин на 4 эпохи у аналога, представление на разрешении $512^3$ компактнее плотной сетки в 4,8--15 раз, пиковая видеопамять на этом разрешении -- 3,3--3,4~ГиБ против 10,4~ГиБ (п.~4.1.23 ТЗ);
    \item метод реконструкции освещения (п.~3.1.9 ТЗ, раздел~\ref{chap:light-comparison}): PSNR изображения сцены с восстановленным источником -- 38,4--50,4~дБ при требовании 25~дБ (п.~4.1.17 ТЗ); время вычисления градиента -- 5,5--11,6~мс на графическом процессоре RTX~2080 и 11--29~мс на центральном процессоре при пределе 10 минут (п.~4.1.18 ТЗ); одинаковый с системой Mitsuba~3 уровень PSNR достигается в 2,6--4,3 раза быстрее на том же графическом процессоре, а итоговая ошибка восстановленной текстуры в 2,3--7,6 раза меньше (п.~4.1.19 ТЗ);
    \item система LightCadRenderCore (п.~3.1.13 ТЗ, подраздел~\ref{sec:lightcad-tests}): при входном изображении $1024\times1024$ время вычисления градиента на графическом процессоре RTX~2080 составляет 6,0--11,4~мс на шаг при пределе 10 минут (п.~4.1.24 ТЗ).
\end{itemize}

Публикации (п.~3.2 ТЗ): по результатам работы опубликована статья Budak~A., Garifullin~A., Frolov~V. SparseDR: Differentiable Rendering of Sparse Signed Distance Fields в трудах конференции IJCAI-26 (уровень A*) \todo{статьи в журналы Q1}.

В рамках данной работы получены новые оригинальные результаты, которые планируется развивать в нескольких направлениях: учёт вторичного освещения при реконструкции поверхностей, совместная реконструкция геометрии, материалов и освещения, реконструкция источников по нескольким изображениям и по фотографиям реальных светильников.

\endinput
